\section{Discussion and limitations}\label{sec:discussion}

\subsection{Interpretation}
The primary result is a solved backward channel at a free plate edge in a
piezoelectric-semiconductor plate with the retained carrier--strain-gradient
coupling. The edge residual and control-volume power balance are at the level of
machine roundoff, so the edge matching and power accounting are numerically
consistent. The isolated backward self-power at \SI{2}{\micro\metre} is
\SI{2.315}{\percent} of the incident self-power for \(H=0\). This is a local
self-flux ratio at one reference plane, not total reflectance and not a
broad-angle experimental result.

The \(H\)-normalized cases show a small, monotonic change in the isolated
backward self-power: \SI{2.286}{\percent} at \(H/H^*=-0.05\), \SI{2.315}{\percent}
at \(H/H^*=0\), and \SI{2.266}{\percent} at \(H/H^*=+0.05\). The sign of the
change follows the stored backward self-flux data. No physical calibration of
\(H\) is claimed.

\subsection{What is validated}
The retained tensor sectors have the expected \textit{6mm} minor and major
symmetries. The \(A\) sector is constructed from \(C\otimes I\) and has the
required major symmetry. The retained \(H\) direction satisfies the strict full-3D
Schur-stability bound used in the production script. The classical
\(\ell=H=0\) control reproduces the independent eight-state classical solver within
the stored tolerance. The edge matching uses SVD and reports rank deficiency as a
failure.

\subsection{What is not validated}
The open-boundary electrical runs are singular in all tested cases. They are
reported as failures, not as successful comparisons. The Qu equivalence and novelty
question remains unresolved because the primary Appendix C equations were not
accessible. No first-ever claim is made.

The synthetic \(A\), \(f\), and \(H\) values are not calibrated ZnO constants. The
retained carrier pair is a primary bulk Hall transfer with stated Hall-to-drift
and thinning assumptions; the diffusivity is derived from the Einstein relation.
The base elastic, piezoelectric, dielectric, and carrier constants are literature-model
values, not a joint same-specimen uncertainty analysis.

\subsection{Mathematical and numerical limitations}
The three-displacement calculation is a 2.5D Fourier-reduced geometry with a
propagation-plane tangential component. It is not a strict sagittal plane-strain
calculation. No full arbitrary 3D edge solution is claimed. Finite-distance H-return
changes are reported as isolated self-flux transfer and incident self-flux factors;
edge conversion and propagation attenuation are not generally separable at finite
distance. The energy velocity is computed from stored \(P_x/W\) values; the earlier
\(|\omega|/|k|\) representation was corrected.

\subsection{Remaining research gaps}
A complete negative-reflection study would require a measured or independently
derived \(H\) tensor, a full calibrated flexoelectric tensor, an open-boundary
solver that avoids the singular edge system, a broad-angle and multi-frequency
experimental comparison, and a primary-source comparison with Qu and co-workers.
These are future work, not established here.
